\section*{Remerciements}
\addcontentsline{toc}{section}{Remerciements}

La réalisation de ce mémoire de fin d'études n'aurait pas été possible sans
le concours de nombreuses personnes à qui je tiens à exprimer ma profonde
gratitude.

Mes remerciements les plus sincères vont en premier lieu à mon directeur de
mémoire, Dr Mamadou Abdoulaye \textsc{Diallo}, dont les orientations
rigoureuses, la disponibilité et les conseils éclairés ont guidé chaque étape
de ce travail. Sa rigueur scientifique et son exigence intellectuelle ont été
pour moi une source d'émulation constante.

J'exprime ma gratitude à Monsieur le Directeur général de l'Agence nationale
de la statistique et de la démographie, dont l'institution produit les enquêtes
sur lesquelles repose ce travail, ainsi qu'à Monsieur Idrissa \textsc{Ndiaye}, Directeur de l'École nationale de
la statistique et de l'analyse économique, pour les conditions de formation et
d'encadrement offertes tout au long de ce cycle.

J'adresse également mes remerciements à Monsieur Omar \textsc{Sène} pour ses
enseignements en évaluation d'impact.

Ma gratitude va aussi à Monsieur Thierno Ibrahima \textsc{Barry}, responsable
de la filière Ingénieur statisticien économiste, et à Dr Mamadou
\textsc{Baldé}, responsable des études et des stages, pour leur accompagnement
tout au long de ce cycle.

Je remercie l'ensemble du corps enseignant de l'École nationale de la
statistique et de l'analyse économique (ENSAE) pour la qualité
de la formation dispensée tout au long de ces années. Les enseignements reçus
ont constitué le socle sur lequel repose la démarche méthodologique de cette
étude.

Je tiens à exprimer ma reconnaissance aux aînés des promotions précédentes de
la filière Ingénieur statisticien économiste, pour leurs conseils, le partage
de leur expérience du mémoire et l'aide qu'ils m'ont apportée aux moments les
plus difficiles de ce travail.

Ma reconnaissance va aussi à mes camarades de classe et de
promotion, pour les échanges fructueux, les moments de solidarité et
l'ambiance de travail stimulante qui ont marqué ces années de formation à
l'ENSAE.

Enfin, je rends hommage à ma famille, et plus particulièrement à mes parents,
dont le soutien moral et les sacrifices consentis ont été le moteur de ma
persévérance. Ce travail leur est dédié.

Que tous ceux qui, de près ou de loin, ont contribué à l'aboutissement de
ce mémoire trouvent ici l'expression de ma sincère reconnaissance.
